% Complete Section 8.1, original PDF 107-108 / printed 97-98.
% block: eigenvalue-problems-heading page:107 type:section
\subsection{{特征值问题}}\label{sec:eigenvalue-problems}

% block: hodge-eigenvalue-problem page:107 type:ordinary-paragraph
{除 Hodge Laplacian 的源问题外, 我们常常还希望求解相应的特征值问题}
\[
(\delta d+d\delta)u=\lambda u,
\]
{连同边界条件, 其中 $\lambda\in\mathbb R$, $u\neq0$. (注意, $\lambda=0$ 是特征值当且仅当存在非零调和形式, 并且调和形式构成零特征值对应的特征空间.)}

% block: hodge-eigenvalue-mixed-formulation page:107 type:ordinary-paragraph
{令 $\sigma=d^*u$, 得到特征值问题的混合弱形式: 求 $\lambda\in\mathbb R$, $\sigma\in H\Lambda^{k-1}$, $0\neq u\in H\Lambda^k$, 使得}
\[
\begin{aligned}
\langle\sigma,\tau\rangle-\langle u,d\tau\rangle&=0,&\tau&\in H\Lambda^{k-1},\\
\langle d\sigma,v\rangle+\langle du,dv\rangle&=\lambda\langle u,v\rangle,&v&\in H\Lambda^k.
\end{aligned}
\]

% block: hodge-eigenvalue-finite-element-discretization pages:107,108 type:ordinary-paragraph
{此时有限元离散化就很清楚了. 将空间 $H\Lambda^{k-1}$ 和 $H\Lambda^k$ 替换为针对源问题已经导出的稳定有限元微分形式空间对 $V_h^{k-1}$ 和 $V_h^k$. 这就得到对称稀疏广义矩阵特征值问题, 可用 Krylov–Schur 方法等各种特征值迭代求解器求解.}

% block: eigenvalue-convergence-requirements page:108 type:ordinary-paragraph
{特征值问题的收敛理论相当精细. 一个好的收敛结果当然必须说明, 若网格充分细, 连续问题 (无穷多个) 特征值中的任意一个, 都能由离散问题的某个特征值任意精确地逼近. 还必须考虑重数, 使连续问题的一个 $m$ 重特征值, 由离散问题的 $m$ 个特征值逼近 (计入它们的重数, 该重数通常不同于 $m$). 还必须排除伪特征值, 即证明随着网格加密, 离散特征值不会在任何并非连续问题特征值的数处聚积. 对重特征值, 尤其必须谨慎陈述特征向量的逼近.}

% block: eigenvalue-convergence-analysis page:108 type:ordinary-paragraph
{幸运的是, 特征值问题的 Galerkin 逼近收敛理论已得到细致发展, 并适用于 Hodge Laplacian. 基本理论由 Babuška 和 Osborn \BookCite{16} 建立. \BookCite{18} 给出了出色的阐述. [\BookNumber{11}, Section 8] 和 [\BookNumber{13}, Section 3.6] 简要讨论了它在 Hodge Laplacian 特征值问题中的应用, 这里不再重复. 主要结论是, Babuška–Osborn 理论的假设可由 Theorem~\ref{thm:mixed-improved-errors} 的改进误差估计推出. 因此, 只要该定理的假设, 即 Hilbert 复形的紧性性质以及 Galerkin 空间子复形的 $W$-有界上链投影得到满足, 就有特征值和特征向量的收敛性. 特征向量的收敛阶是最优的 (由空间的逼近性质决定), 而特征值的收敛阶是特征向量收敛阶的两倍.}

% block: hodge-and-maxwell-spectral-examples page:108 type:ordinary-paragraph
{在引言的 Table~\ref{tab:feec-vector-eigenvalues} 中, 已经看到了 Hodge Laplacian \eqref{eq:vector-laplacian-eigenproblem} 特征值收敛的一个例子. 与此相反, Table~\ref{tab:vector-lagrange-wrong-eigenvalues} 表明, 一个不基于混合形式的朴素方法并不收敛. 引言中还考虑了 Maxwell 特征值问题 $\operatorname{curl}\operatorname{curl}u=\lambda u$. 基于 Lagrange 单元的朴素方法导致错误结果, 真正的谱或者无法从计算中辨认, 或者受到持续存在的伪特征值污染, 如 Table~\ref{tab:maxwell-spectral-comparison} 所示. 使用弱形式和最低阶 Whitney 1-形式进行的计算, 在表中标为 FEEC 的各行中给出, 收敛良好. 将 Maxwell 特征值问题与 Hodge Laplacian 特征值问题联系起来就可证明这一点. 前者在 Hodge 分解后作为 $\mathcal B_1^*$ 问题出现, 因而 Hodge Laplacian 特征值问题的收敛分析可推广到 Maxwell 问题. 细节见 [\BookNumber{13}, Section 3.6.1].}
